% sections/network_construction.tex --- Network Construction
\section{Network Construction}
\label{sec:network}

This section makes precise the mapping from the row-level
journal-entry table to a directed account-level graph.  The
treatment follows the formulation in
\citet{ivertowski2024apn} but is presented self-contained.

\subsection{Setup and notation}
\label{sec:network:setup}

Let $\mathcal{A}$ be the set of general-ledger accounts in the
chart of accounts.  A \emph{journal entry} $J$ is a balanced pair
$(D,C)$ of multisets of line items
\[
  D = \{(a^D_1, d_1), \ldots, (a^D_U, d_U)\},
  \quad
  C = \{(a^C_1, c_1), \ldots, (a^C_V, c_V)\},
\]
with $a^D_u, a^C_v \in \mathcal{A}$ and $d_u, c_v \in \R_{>0}$
satisfying the balance constraint
\begin{equation}
  \sum_{u=1}^{U} d_u \;=\; \sum_{v=1}^{V} c_v. \label{eq:balance}
\end{equation}
We refer to $U = |D|$ and $V = |C|$ as the debit and credit
\emph{cardinalities}.  Bin~1 of
Section~\ref{sec:empirical:line-items} corresponds to
$U + V \in \{2, \ldots, 9\}$ and accounts for
$92.60\%$ of empirical postings.

The \emph{accounting graph} associated with a collection of journal
entries over a period $[t_0, t_1]$ is the directed multigraph
\begin{equation}
  G(t_0, t_1) \;=\; \big( A(t_0, t_1),\, E(t_0, t_1) \big),
  \label{eq:graph}
\end{equation}
where the node set $A(t_0, t_1) \subseteq \mathcal{A}$ contains
every account that carries a non-zero balance over the period and
the edge set $E(t_0, t_1)$ contains a weighted, signed,
time-stamped edge for every $(a^D, a^C)$ pair generated by some
journal entry posted in $[t_0, t_1]$.  When edges between the same
ordered account pair are aggregated, the period-level edge weight is
\begin{equation}
  E_{(a,b)}(t_0, t_1) \;=\; \sum_{n=1}^{|\Omega_{(a,b),[t_0,t_1]}|}
                                \omega_n, \label{eq:edge-aggregation}
\end{equation}
where $\Omega_{(a,b),[t_0,t_1]}$ is the multiset of individual
transaction flows from $a$ to $b$ posted in the period and
$\omega_n$ is the monetary amount of the $n$-th such flow.  Equation
\eqref{eq:edge-aggregation} reproduces \citet{ivertowski2024apn}
Eq.~(2).

The construction of the per-entry edge set is the central object of
this section.  Five methods (A through~E) are defined; the choice
of method depends on the cardinalities $(U, V)$ and on whether the
line-item amounts are pairwise distinct.  Methods are tried in
order; the first that succeeds is the one that materialises the edge
set for the entry.  Algorithm~\ref{alg:network-gen} states the
dispatch.

\subsection{Method A: 1-to-1 mapping}
\label{sec:network:A}

Method~A applies when exactly one debit and one credit line
($U = V = 1$, hence $U + V = 2$).  By the balance constraint
\eqref{eq:balance} the single debit amount equals the single credit
amount, and a single edge $(a^D_1, a^C_1)$ with weight $d_1 = c_1$
is materialised:
\[
  E_J^A \;=\; \{\, (a^D_1, a^C_1, d_1) \,\}.
\]
Method~A is exact: every line item maps bijectively to the unique
debit-credit pair, and the \emph{confidence} of the inferred edge
is $1$ (Table~\ref{tab:confidence}).  Per
Section~\ref{sec:empirical:line-items}, Method~A applies to
$60.68\%$ of postings in the corpus.

\subsection{Method B: n-to-n mapping}
\label{sec:network:B}

Method~B applies when $U = V \geq 2$, that is, when an equal
number of debit and credit line items appear.  Method~B takes the
Cartesian product of debit and credit lines:
\[
  E_J^B \;=\; \big\{\, (a^D_u, a^C_v, w_{uv}) :
    u \in \{1,\ldots,U\}, v \in \{1,\ldots,V\} \,\big\},
\]
where the weight $w_{uv}$ is $\min(d_u, c_v)$ when the line-item
amounts are pairwise distinct (so the matching is forced), and
$\min(d_u, c_v)$ scaled by a confidence factor when they are not.
The confidence assigned to a non-distinct edge is
$1/n_i$ where $n_i$ is the number of indistinguishable values the
$i$-th line item takes (Table~\ref{tab:confidence}).  Method~B
preserves total monetary flow ($\sum_{u,v} w_{uv}$ equals the
balanced amount) but does not in general preserve the per-line
allocation when amounts collide.

\subsection{Method C: n-to-m mapping}
\label{sec:network:C}

Method~C applies when $U \neq V$ and the entry's line-item amounts
admit a partition that satisfies the balance constraint
on both sides.  When $V > U$ (more credit lines than debit
lines), this is the classical problem of partitioning the $V$
credit amounts into $U$ groups whose sums match the $U$ debit
amounts.  The number of admissible partitions is
\begin{equation}
  a(n, m) \;=\; \frac{n!}{m!} \sum_{j=0}^{m}
                (-1)^{m-j} \binom{m}{j} j^n, \label{eq:partition-count}
\end{equation}
reproducing Eq.~(3) of \citet{ivertowski2024apn} with $n = V$ and
$m = U$.  When this combinatorial count exceeds a configurable
budget---which it does rapidly for large $n, m$---Method~C is
abandoned in favour of Methods~D or~E.

When the partition exists and is admissible, Method~C materialises
one edge per (debit-line, credit-group) pair with weight equal to
the group sum.  Confidence is $1$ when the partition is unique and
$1/n_i$ otherwise (Table~\ref{tab:confidence}).

\subsection{Method D: composition to higher aggregate}
\label{sec:network:D}

Method~D handles the case where Methods~B and~C fail by
\emph{composing} line items up to a higher aggregate of the chart
of accounts (account sub-class \texttt{ASC} or account class
\texttt{AC}).  Aggregation is exact---debits and credits sum
unchanged---and may collapse the cardinality mismatch
($U \neq V$) into the equal case ($U = V$), in which event
Method~B can then be applied at the aggregated level.  Method~D
re-runs the dispatch chain (Methods~A, B, C, in that order) on the
aggregated entry; if it succeeds, the resulting edges connect
\emph{aggregate} nodes rather than individual accounts.  Confidence
is~$1$ when the aggregated entry admits an exact lower-level
solution (Table~\ref{tab:confidence}).

\subsection{Method E: decomposition with shadow line items}
\label{sec:network:E}

Method~E is the residual fallback when no exact solution exists at
any aggregation level.  It decomposes credit-side amounts into
shadow line items that match debit-side amounts, materialising
edges with confidence $1/n^b$ where $n^b$ is the number of
decomposition steps.  Method~E always produces a result and so
guarantees that the algorithm terminates with a valid edge set; it
is, however, the lowest-confidence method, and its prevalence on a
particular dataset is itself a quality signal.

\subsection{Confidence}
\label{sec:network:confidence}

\begin{table}[H]
  \centering
  \caption{Per-method edge confidence.  Source: \citet{ivertowski2024apn},
    Table~V.  $n_i$ denotes the count of the $i$-th line-item value;
    $n^b$ denotes the count of decomposition steps.}
  \label{tab:confidence}
  \begin{tabular}{llc}
    \toprule
    \textbf{Method} & \textbf{Case} & \textbf{Confidence}\\
    \midrule
    A & ($U = V = 1$, distinct)            & $1.00$\\
    B & ($U = V \geq 2$, distinct)         & $1.00$\\
    B & ($U = V \geq 2$, not distinct)     & $1.00 / n_i$\\
    C & ($V > U$, distinct)                & $1.00$\\
    C & ($V > U$, not distinct)            & $1.00 / n_i$\\
    C & ($V < U$)                          & $1.00 / n^b$\\
    D & (any case, exact aggregation)      & $1.00$\\
    E & (decomposition fallback)           & $1.00 / n^b$\\
    \bottomrule
  \end{tabular}
\end{table}

The journal-entry-level confidence is the product of the line-level
confidences over all edges produced for the entry, consistent with
the treatment in \citet{ivertowski2024apn}.  This per-entry
confidence is exposed on every row of \texttt{graphs/je\_network.csv}
in the \texttt{confidence} column (Section~\ref{sec:graph}).

\subsection{Dispatch}
\label{sec:network:dispatch}

\begin{algorithm}[H]
  \DontPrintSemicolon
  \caption{Network generation (after \citet{ivertowski2024apn},
    Algorithm~1).}
  \label{alg:network-gen}
  \KwIn{Journal entries $J_1, \ldots, J_N$; chart of accounts
    aggregator $\mathcal{M}$.}
  \KwOut{Edge sets $\mathrm{NR}_1, \ldots, \mathrm{NR}_N$.}
  \For{$n \leftarrow 1$ \KwTo $N$}{
    $U \leftarrow |D(J_n)|$;
      $V \leftarrow |C(J_n)|$\;
    \uIf{$U + V = 2$}{
      $\mathrm{NR}_n \leftarrow \mathrm{Method\,A}(J_n)$
    }
    \uElseIf{$U = V$}{
      $\mathrm{NR}_n \leftarrow \mathrm{Method\,B}(J_n)$
    }
    \uElseIf{$\mathrm{Method\,C}(J_n)$ \emph{is feasible}}{
      $\mathrm{NR}_n \leftarrow \mathrm{Method\,C}(J_n)$
    }
    \Else{
      $J^{\mathrm{ASC}}_n \leftarrow
        \mathrm{Method\,D}(J_n, \mathcal{M}, \texttt{ASC})$\;
      \uIf{dispatch on $J^{\mathrm{ASC}}_n$ succeeds}{
        $\mathrm{NR}_n \leftarrow$ dispatch result
      }
      \uElse{
        $J^{\mathrm{AC}}_n \leftarrow
          \mathrm{Method\,D}(J_n, \mathcal{M}, \texttt{AC})$\;
        \uIf{dispatch on $J^{\mathrm{AC}}_n$ succeeds}{
          $\mathrm{NR}_n \leftarrow$ dispatch result
        }
        \uElse{
          $\mathrm{NR}_n \leftarrow \mathrm{Method\,E}(J_n)$
        }
      }
    }
  }
  \KwRet{$\mathrm{NR}_1, \ldots, \mathrm{NR}_N$}
\end{algorithm}

\subsection{Default in the released benchmark}
\label{sec:network:default}

For \bench{} the default construction is Method~A on $1$-to-$1$
entries and Method~B on $n$-to-$n$ entries with $U = V$.  Methods
C--E remain available for ablation studies where a reader wants to
test sensitivity to construction strategy; they are activated by the
\texttt{network\_construction} option in the released configuration
(Appendix~\ref{sec:appendix:config}).  This default reproduces the
edge counts reported in Section~\ref{sec:evaluation}; alternative
configurations produce different counts by construction and hence
require a re-baselining of all reported numbers.

\subsection{Equivalence at the period level}
\label{sec:network:equivalence}

Although Methods~A through~E differ in how they distribute monetary
flow across edges \emph{within} a single posting, they are
\emph{equivalent at the period level} in the sense of
Equation~\eqref{eq:edge-aggregation}: for any two methods $M$ and
$M'$ in the family, the aggregated edge weight
$E_{(a,b)}(t_0, t_1)$ is the same when summed over all postings in
the period.  This follows from the balance constraint
\eqref{eq:balance}: the total monetary flow into and out of any
account is determined by the line items of the entries, not by
the way those line items are paired.  The methods therefore
correspond to different \emph{micro-level} resolutions of the same
\emph{macro-level} consolidation, in direct analogy to the
financial-statement consolidation operation in which subsidiary
ledgers are aggregated to group ledgers.
